\section{滑模方法2（积分滑模）}

定义 $\dot\eta_i=k_1p_{i0}-\varphi_i$，滑模变量 $s_{i0}=\eta_i+k_2p_{i0}+v_{i0}$，其中 $\varphi_i$ 同方法1。由 $\dot p_{i0}=v_{i0}$、$\dot v_{i0}=B_iu_i+w_i$ 与 $\dot\eta_i$ 得
\[
    \dot s_{i0}=B_i(t)u_i+\Delta_i+w_i,\qquad
    \Delta_i=k_1p_{i0}-\varphi_i+k_2v_{i0}.
\]
逐分量 $\dot s_{i0,\ell}=b_{i\ell}u_{i\ell}+N_{i\ell}$，控制律与增益同方法1。注意 $\Delta_i$ 只含 $\varphi_i$（位置量），无需 $\dot\varphi_i$，故不依赖邻居速度。

\subsection{到达性}

与方法1 完全相同：关键不等式 $|b|R>|N|$ 仍成立，故各分量有限时间到达相应周期面，$s_{i0}\in L_\infty$，滑模面上 $\dot s_{i0}=0$。$b_{i\ell}$ 有限次切换时最后一次后进入滑模；无限切换需驻留时间假设。

\subsection{降阶名义合围动力学}

由 $\dot s_{i0}=0$ 与 $v_{i0}=\dot p_{i0}$ 得
\[
    \ddot p_{i0}+k_2\dot p_{i0}+k_1p_{i0}=\varphi_i,
\]
即标准的“目标吸引 $k_1p_{i0}$+阻尼 $k_2\dot p_{i0}$+智能体间排斥 $\varphi_i$”名义合围系统。该形式与已有合围编队文献完全一致，故凸包合围、碰撞避免与速度一致性可由已有名义系统结论直接推出。

说明：本方法把 $\varphi_i$ 关入积分器，既无需邻居速度，又得到标准降阶动力学，可直接引用文献结论。技术条件包括：输入增益须为对角结构、控制方向不穿越零、切换不能无限快、滑模意义为理想/Filippov、合围结论引用名义系统文献。
